% Partie : set-theory
% Chapitre : ordinals
% Section : ordtype

\documentclass[../../../include/open-logic-section]{subfiles}

\begin{document}

\olfileid[fr]{sth}{ordinals}{ordtype}
\olsection{Les ordinaux comme types d'ordre}

Munis du remplacement, et travaillant donc à présent dans
$\ZFminus$, nous pouvons enfin démontrer le résultat recherché :

\begin{thm}\ollabel{thmOrdinalRepresentation}
Tout bon ordre est isomorphe à un unique ordinal.
\end{thm}

\begin{proof}
Soit $\tuple{A, <}$ un bon ordre. Par la
\olref[basic]{ordisoidentity}, il est isomorphe à au plus un ordinal.
Pour raisonner par l'absurde, supposons qu'il ne soit isomorphe à
\emph{aucun} ordinal. Nous allons d'abord rendre $\tuple{A, <}$
«~aussi petit que possible~». Plus précisément, s'il existe un
segment initial propre $\tuple{A_a, <_a}$ qui n'est isomorphe à
aucun ordinal, choisissons le plus petit $a \in A$ qui possède cette
propriété, puis posons $B = A_a$ et $\mathord{\lessdot} =
\mathord{<_a}$. Sinon, posons $B = A$ et $\mathord{\lessdot} =
\mathord{<}$.

Par construction, chaque segment initial propre de $B$ est isomorphe
à un ordinal, unique d'après ce qui précède. Le remplacement assure
donc l'existence de l'ensemble suivant, qui est une fonction :
\[
	f = \Setabs{\tuple{\beta, b}}{b \in B\text{ et }\ordeq{\beta}{\tuple{B_b, \lessdot_b}}}
\]
Pour achever la contradiction, nous allons montrer que $f$ est un
isomorphisme $\alpha \to B$ pour un certain ordinal $\alpha$.

Il est clair que $\ran{f} = B$. Par le
\olref[iso]{lemordsegments}, $f$ préserve l'ordre, c'est-à-dire que
$\gamma \in \beta$ si et seulement si
$f(\gamma) \lessdot f(\beta)$. Pour montrer que $\dom{f}$ est un
ordinal, il suffit, par le \olref[basic]{corordtransitiveord}, de
montrer qu'il est transitif. Fixons donc $\beta \in \dom{f}$ : pour
un certain $b$, nous avons
$\ordeq{\beta}{\tuple{B_b, \lessdot_b}}$. Si
$\gamma \in \beta$, alors $\gamma \in \dom{f}$ par le
\olref[iso]{wellordinitialsegment}. Ainsi,
$\beta \subseteq \dom{f}$.
\end{proof}

Ce résultat autorise la définition suivante, que nous souhaitions
donner depuis la \olref[vn]{sec} :

\begin{defn}
Si $\tuple{A, <}$ est un bon ordre, son type d'ordre,
$\ordtype{A, <}$, est l'unique ordinal $\alpha$ tel que
$\ordeq{\tuple{A, < }}{\alpha}$.
\end{defn}

Cette définition permet aussi d'établir deux propriétés utiles :

\begin{cor}\ollabel{ordtypesworklikeyouwant}
Si $\tuple{A, <}$ et $\tuple{B, \lessdot}$ sont des bons ordres :
\begin{align*}
	\ordtype{A, <} = \ordtype{B, \lessdot}&\text{ ssi }\ordeq{\tuple{A, <}}{\tuple{B, \lessdot}}\\
	\ordtype{A, <} \in \ordtype{B, \lessdot}&\text{ ssi }\ordeq{\tuple{A, <}}{\tuple{B_b, \lessdot_b}}\text{ pour un certain }b \in B
\end{align*}
\end{cor}

\begin{proof}
La première égalité découle de la \olref[basic]{ordisoidentity}. Pour
démontrer la seconde, posons $\ordtype{A, <} = \alpha$ et
$\ordtype{B, \lessdot} = \beta$, et soit $f \colon \beta \to B$
notre isomorphisme d'ordre.
Si $\alpha \in \beta$, la restriction de $f$ à $\alpha$ est un
isomorphisme de $\alpha$ sur $B_{f(\alpha)}$. Ainsi,
$\tuple{A, <}$ est isomorphe à ce segment initial de $B$.
Réciproquement, supposons $\tuple{A, <}$ isomorphe à $B_b$ pour
un certain $b \in B$. Posons $\gamma=f^{-1}(b)$ ; alors
$\gamma \in \beta$ et la restriction de $f$ à $\gamma$ est un
isomorphisme sur $B_b$. Les ordinaux $\alpha$ et $\gamma$ sont
donc isomorphes, puis égaux. Ainsi, $\alpha \in \beta$.
Ces arguments s'appuient sur les résultats donnés
par la \olref[iso]{ordisounique},
le \olref[iso]{wellordinitialsegment} et le
\olref[basic]{ordissetofsmallerord}.
\end{proof}

\end{document}
